An intelligent control system, such as a mobile robot, may be required to perform many different tasks.
If the \envword\ set is finite, like selecting between ``map an environment'' and ``deliver a package'',
then its size is naturally quantified by the number of \envwords.
If the \envword\ set is infinite, like delivering packages with arbitrary mass and inertial properties,
then its size is not so easily quantified.
Even in a metric or measure space,
these structures may be only weakly linked to the amount of variation in behavior
required for good performance on all \envwords.

Our interest in this issue is motivated by \multienv\ paradigms in learning-based control,
where the policy is selected from a parameterized family of functions
that map state and \envword\ parameters directly to actions.
As the \envword\ space expands from a singleton set,
we expect to need a more expressive class of functions
to represent a good \multienv\ policy.
In this work, we propose the 
\emph{$\alpha$-suboptimal covering number}
to capture this idea. %
For a \envword\ space $\Envs$
and a suboptimality ratio $\alpha > 1$,
we define $\coveringnum{\alpha}{\Envs}$ as
the size of the smallest set of \singleenv\ policies $\coveringset$ %
such that for every $\env \in \Envs$, at least one $\pi \in \coveringset$
has a cost ratio no greater than $\alpha$ relative to the optimal policy for $\env$.
If the policies in $\coveringset$ are parameterized functions,
then $\coveringset$ provides an upper bound on the number of parameters
needed to represent an $\alpha$-suboptimal \multienv\ policy.

To study suboptimal covering numbers in a concrete setting,
we consider linear dynamical systems
with quadratic cost functions, or LQR problems.
We construct a family of well-behaved \multienv\ LQR problems 
where $\Envs$ is controlled by a ``breadth'' parameter $\breadth \in [1, \infty)$,
and for which
$\coveringnum{\alpha}{\Envs_\breadth}$
is always finite but increasing in $\breadth$.
LQR problems are a common setting to analyze learning algorithms because detailed properties are known \citep[for example]{fazel-global-lqr}.
This has led to new inquiries into their fundamental properties \citep{bu-topological-gains}.

For the special case of a scalar LQR problem,
we derive matching logarithmic upper and lower bounds on
$\coveringnum{\alpha}{\Envs_\breadth}$
as a function of $\breadth$.
As an effort towards analogous bounds for the matrix case,
we present empirical results
and collect some useful tools from the control literature.
For upper bounding, we analyze properties of a logical extension of our scalar cover into a geometric grid.
For lower bounding, we compare two choices of ``extremal''  \multienv\ LQR problems:
minimal coupling ($A$ is an identity matrix),
and maximal coupling ($A$ is a scaled matrix of all ones).
For minimal coupling, we find surprising topological behavior of the suboptimal neighborhoods.

This paper is an initial step towards a comprehensive theory.
In addition to a more complete picture of deterministic LQR systems,
ideas of $\alpha$-suboptimal covers could be applied to a wide range of \multienv\ problems.
We also hope they will lead to insights
about function class expressiveness in learning-based \multienv\ control.



